F4PGA brands itself as the GCC of FPGAs. 
While most FPGA architectures are not open source, F4PGA has functional flows for both the few more open companies like Lattice 
and more more closed companies like Xilinx.
Project x-ray is an attempt to reverse engineer the Xilinx architectures from compiled bitstreams to create funcitonal open source toolchains.
Unfortunately at time of writing F4PGA is not pre-compiled for ARM silicon; however, this is luckily not a deal-breaker for us.

In theory we could recompile everything included in F4PGA; however, 
the instructions for F4PGA rely heavily on everything being pre-packaged, 
and the individual tools typically assume that you are operating on an x86 System. 
In order to get this working you are going to need to use Apple's Rosetta for x86 emulation, 
and docker to containerize the needed software. 
Previous instructions have pointed towards using a Windows VM and running Vivado within that VM; 
while this works, it is less than ideal due to the high overhead to run Windows within OSx, 
especially with the additional overhead of emulation. 
To this end, these instructions will take the following alternate approach.

\begin{itemize}
    \item Install Docker. 
    \item Pull or Build the needed Docker Image.
    \item Run the docker image with reference code to produce a bitstream.
    \item Load that bitstream over to an attached Basys 3 FPGA board.
\end{itemize}

\subsubsection{Install Docker}
To install docker on Apple Silicon, you are going to want Docker Desktop and for our purposes ensure that Rosetta 2 is installed and enabled.

To install Rosetta run the following: 

\begin{lstlisting}[language=bash]
    softwareupdate --install-rosetta
\end{lstlisting}

To install Docker Desktop on apple silicon go to the following link: 

\begin{lstlisting}[language=bash]
    https://desktop.docker.com/mac/main/arm64/Docker.dmg?utm_source=docker&utm_medium=webreferral&utm_campaign=docs-driven-download-mac-arm64
\end{lstlisting}

Run the installer and let the Docker Daemon launch.

\subsubsection{Pull the Needed Image}
With docker instaled you now need an image to run.

To pull the image:
\begin{lstlisting}[language=bash]
    docker pull ghcr.io/devinatkin/MAC-FPGA-SETUP:latest
\end{lstlisting}

If you would prefer to build the image yourself, go to the repo for this document and grab the reference docker file. 
With that saved in the local directory the following can be run to build it:

\begin{lstlisting}[language=bash]
    docker build --platform=linux/amd64 -t f4pga-docker .
\end{lstlisting}

This builds the docker container and specifies to build it as an x86 image. 
\textit{This may seem strange given you are running on an ARM device; however, you will need to use Rosetta as not all dependencies are available on apple silicon.}


\textbf{Be warned that the images for this are very large (~20GB) and will take awhile to download or build.}


\subsubsection{Produce test bitstream}
You have now installed docker and pulled the required image. 
The next step is to run the image and produce a test bitstream to verify that everything is working correctly.

\begin{lstlisting}[language=bash]
    docker run --rm -it --platform linux/amd64 f4pga-docker /bin/bash
\end{lstlisting}

This will run the f4pga docker container and provide an interactive prompt. 
This can be a good way to verify that everything is correctly installed. 
\textit{It is very important that the platform is specified, otherwise docker will attempt to run utilizing the native ARM architecture}. 
Once you are done interacting with the docker container run exit to stop and remove the container.

\begin{lstlisting}[language=bash, columns=fullflexible]
git clone https://github.com/chipsalliance/f4pga-examples

docker run --rm -it --platform linux/amd64 \
    -v /Users/your_username/f4pga/f4pga-examples/xc7:/opt/f4pga-examples/xc7 \
    -w /opt/f4pga-examples/xc7 \
    f4pga-docker /bin/bash -i -c "make -C counter_test TARGET='basys3'"
\end{lstlisting}

Ideally we'll actually want to run a bitstream generation and get the results outside of the docker container. 
For some test code we can clone the F4PGA-Examples repository and use the provided example code so that we can reach a bitstream as quickly as possible. 
Run the above commands in your terminal, ensuring to appropriately substitute out your own directories. 

You'll notice in the docker run that a few additional elements have been added to the commands.
\begin{itemize}
    \item --rm option, this removes the container once execution stops. 
    This means that you don't need to manually remove it before executing again.

    \item --it option, this pretties up the output and directs the output to the main terminal. 
    It will run without this; however, you won't know the current state or if anything has gone wrong. 

    \item -v option, this is the volume mount. I'm pointing this towards the directory of the f4pga example code targeting the Artix 7 series. 
    You'll want to adjust the directory to match your own directory here. 
    This also means that once the bitstream is generated you'll see it appear in your native file system without needing to transfer the data manually. 

    \item -w option, this option you may or may not need depending on your directory structure. It is setting the working directory within the docker container. 

    \item f4pga-docker, this is the container name. Yours may be called MAC-FPGA-SETUP or something else depending on how you pulled the image.

    \item \texttt{/bin/bash -i -c "make -C counter\_test TARGET='basys3'"}, 
    This is what actually executes the bitstream synthesis and implementation. We will go over the Makefile structure later in this document.

\end{itemize}

\subsubsection{Load Bitstream over to attached Basys3}

\begin{minted}{bash}
    openFPGALoader -b basys3 your_bitstream.bit
\end{minted}

Not that you have generated your bitstream you can use openFPGALoader which you should have installed earlier using brew. 
Plug the Basys3 board into your Mac and enter in the command as shown, replacing the bitstream file for the bitstream file that you generated previously.
The bitstream should be loaded onto the board and you'll see the ready light on the board turn on. 
Assuming everything went correctly you should not have a programmed FPGA board. 